% ==========================================
% Step 4/4 - Final refinement (refinement step 3)
% Output: step4-10-28-2026-wednesday-week7-Algebra_I-offset-final.tex
% Input: step3-10-28-2026-wednesday-week7-Algebra_I-offset-speech_refined.tex
% Model: gemini-3.8-flash
% Temperature: 0.35
% TopP: 0.9
% TopK: 10
% MaxOutputTokens: 65535
% ThinkingBudget: 4096
% ThinkingLevel: MEDIUM
% Processed on: 2026-10-29 11:13:02
% Thought summaries: step4-10-28-2026-wednesday-week7-Algebra_I-offset-final-thoughts.md
% ------------------------------------------
% Token Usage (1 Request(s)):
%   - Total Prompt Tokens : 61'392 (gesamter Prompt, inkl. Cache)
%   - Cached Context      : 0 (aus Google Context-Cache, rabattiert)
%   - Fresh Input Tokens  : 61'392 (neu verrechneter Input)
%   - Output Tokens       : 27'688 (generiertes LaTeX, ohne Thinking)
%   - Thinking Tokens     : 17'053 (Denk-Tokens, als Output verrechnet)
% ==========================================
% ==========================================
% Step 3/4 - Speech refinement (refinement step 2)
% Output: step3-10-28-2026-wednesday-week7-Algebra_I-offset-speech_refined.tex
% Input: step2-10-28-2026-wednesday-week7-Algebra_I-offset-merged.tex
% Model: gemini-3.8-flash
% Temperature: 0.35
% TopP: 0.9
% TopK: 10
% MaxOutputTokens: 65535
% ThinkingBudget: 4096
% ThinkingLevel: MEDIUM
% Processed on: 2026-10-29 11:00:57
% Thought summaries: step3-10-28-2026-wednesday-week7-Algebra_I-offset-speech_refined-thoughts.md
% ------------------------------------------
% Token Usage (1 Request(s)):
%   - Total Prompt Tokens : 158'466 (gesamter Prompt, inkl. Cache)
%   - Cached Context      : 0 (aus Google Context-Cache, rabattiert)
%   - Fresh Input Tokens  : 158'466 (neu verrechneter Input)
%   - Output Tokens       : 27'528 (generiertes LaTeX, ohne Thinking)
%   - Thinking Tokens     : 11'307 (Denk-Tokens, als Output verrechnet)
% ==========================================
% ==========================================
% Step 2/4 - Merge & timestamp alignment (refinement step 1)
% Output: step2-10-28-2026-wednesday-week7-Algebra_I-offset-merged.tex
% Input: step1-10-28-2026-wednesday-week7-Algebra_I-all-offset.tex
% Model: gemini-3.8-flash
% Temperature: 0.4
% TopP: 0.95
% TopK: 40
% MaxOutputTokens: 65535
% ThinkingBudget: 32765
% ThinkingLevel: MEDIUM
% Processed on: 2026-10-29 10:53:20
% Thought summaries: step2-10-28-2026-wednesday-week7-Algebra_I-offset-merged-thoughts.md
% ------------------------------------------
% Token Usage (1 Request(s)):
%   - Total Prompt Tokens : 117'695 (gesamter Prompt, inkl. Cache)
%   - Cached Context      : 0 (aus Google Context-Cache, rabattiert)
%   - Fresh Input Tokens  : 117'695 (neu verrechneter Input)
%   - Output Tokens       : 28'018 (generiertes LaTeX, ohne Thinking)
%   - Thinking Tokens     : 8'315 (Denk-Tokens, als Output verrechnet)
% ==========================================
% ==========================================
% Step 1/4 - Transcription, 2 part(s) combined
% Output: step1-10-28-2026-wednesday-week7-Algebra_I-all-offset.tex
% Source video: 10-28-2026-wednesday-week7-Algebra_I.mp4
% Model: gemini-3.8-flash
% Temperature: 0.36
% TopP: 0.95
% TopK: 40
% MaxOutputTokens: 65535
% ThinkingBudget: 32765
% ThinkingLevel: HIGH
% Processed on: 2026-10-29 10:44:57
% ------------------------------------------
% Parts:
%   - Part 1: 10-28-2026-wednesday-week7-Algebra_I-part1.mp4 | start 00:00:00 | model gemini-3.8-flash | processed 2026-10-29 09:46:43 | prompt 421'592 | output 14'490 | thinking 16'908 | aus Cache
%   - Part 2: 10-28-2026-wednesday-week7-Algebra_I-part2.mp4 | start 00:19:57 | model gemini-3.8-flash | processed 2026-10-29 10:44:57 | prompt 436'742 | output 17'658 | thinking 39'735
% ------------------------------------------
% Token Usage of this run (1 Request(s)):
%   - Total Prompt Tokens : 436'742 (gesamter Prompt, inkl. Cache)
%   - Cached Context      : 0 (aus Google Context-Cache, rabattiert)
%   - Fresh Input Tokens  : 436'742 (neu verrechneter Input)
%   - Output Tokens       : 17'658 (generiertes LaTeX, ohne Thinking)
%   - Thinking Tokens     : 39'735 (Denk-Tokens, als Output verrechnet)
% ==========================================

% Primary Language: English
% End of the video: 00:43:53

\lecturechapter[7]{Wednesday}{28 Oct}{28 October 2026}{Lagrange's Theorem, Symmetric Groups, and Geometric Symmetries}

\begin{overview}[Course Content Overview]
This lecture explores the foundational structure of finite groups, culminating in Lagrange's Theorem and its geometric applications:
\begin{itemize}
    \item \textbf{Perspectives on Group Theory:} Focusing on finite, non-Abelian groups alongside prototype extreme families: symmetric groups $S_n$ and cyclic groups $\mathbb{Z}_n$.
    \item \textbf{Element Orders and Cyclic Subgroups:} Well-definedness of $\operatorname{Ord}(x)$ and verification that $\langle x \rangle \le G$ is isomorphic to $(\mathbb{Z}_{\operatorname{Ord}(x)}, +, [0])$.
    \item \textbf{Lagrange's Theorem:} Formulation and proof that $|H|$ divides $|G|$ for any subgroup $H \le G$, establishing that element orders divide group orders and that prime-order groups have no non-trivial proper subgroups.
    \item \textbf{Coset Equivalence Relation:} Partitioning $G$ into equivalence classes under $x \sim y \iff \exists h \in H \colon xh = y$ and establishing the bijection $\phi \colon H \to E$, $\phi(h) = xh$.
    \item \textbf{Cycle Decompositions and Partitions:} Dynamic visualization of cycles in $S_n$ and the deep connection between cycle decompositions and integer partitions.
    \item \textbf{Rotational Symmetry Groups of Geometric Solids:} Counting symmetries of the Platonic solids via face stabilization, establishing the geometric isomorphisms $\operatorname{Rot}(\text{Cube}) \cong S_4$ via body diagonals and $\operatorname{Rot}(\text{Tetrahedron}) \cong A_4$ via vertex permutations.
\end{itemize}
\end{overview}

\section{Perspectives on Group Theory}

\subsection{Finite and Non-Abelian Groups}

\begin{speech}[00:00:00 - 00:01:25]
Okay, so, we start. I hope you have now concluded your study of Confucius, Plato, and Schopenhauer, and we can come back to the discussion of groups.

So, I defined last time group --- $G$ is a group \inlinemetanote{writes: $G \text{ group } (G, \cdot, e)$} --- which I don't repeat now. And I remind you, because it's kind of important for the perspective of the discussion about group theory: we will normally... our interest will be here mainly, but not exclusively, focused on finite groups \inlinemetanote{writes: $\implies \text{mostly focused on finite groups}$}, but not always.

And also our main interest, but not always, is on non-commutative groups, so non-Abelian \inlinemetanote{writes: $\implies \text{on non-Abelian}$}. But not always. So, what you should have in your mind when I say group: it's a finite group and not necessarily Abelian.

The two families of examples on extremes are the family of the symmetric groups $S_n$ --- we'll say more; I said a little bit about this before, but I will say more --- symmetric group \inlinemetanote{writes: $S_n \text{ Symmetric group } |S_n| = n!$}. And this is, of course, a finite group of size $n$ factorial. And that's somehow the richest family that we will speak about, and in some sense... in a sense I'll make more precise later, really a very rich family of finite groups, maybe the richest family of finite groups.
\end{speech}

\begin{content}[Standing Focus in Group Theory]
Let $(G, \cdot, e)$ be a group. Throughout this part of the course, our primary analytical focus rests upon:
\begin{itemize}
    \item \textbf{Finite Groups:} $|G| < \infty$ (though infinite groups will occasionally appear).
    \item \textbf{Non-Abelian Groups:} Groups where the operation is generally non-commutative, i.e., $\exists x, y \in G$ with $x \cdot y \neq y \cdot x$.
\end{itemize}

\begin{example}[The Symmetric Group]\label[example]{ex:symmetric-group}
For $n \in \mathbb{N}$, the \newterm{symmetric group} $S_n$ consists of all bijections (permutations) of the set $\{1, 2, \dots, n\}$ under function composition. Its order is given by:
\[
|S_n| = n!
\]
For $n \ge 3$, $S_n$ is non-Abelian.
\end{example}

\begin{steps}[Why Symmetric Groups Are Considered the Richest Family]
The lecturer highlights $S_n$ as \qt{the richest family} of finite groups. Going beyond the lecture: this perspective is formalized by Cayley's Theorem, which establishes that every finite group of order $n$ is isomorphic to a subgroup of the symmetric group $S_n$. Thus, understanding subgroups of symmetric groups in principle encompasses all finite group theory.
\end{steps}
\end{content}

\subsection{Cyclic Abelian Groups}

\begin{speech}[00:01:25 - 00:02:04]
And then there's the simplest family, $\mathbb{Z} / n\mathbb{Z}$ \inlinemetanote{writes: $(\mathbb{Z}_n, +, 0)$}. That's with plus and zero. These are the Abelian groups we've seen before, and these are called cyclic \inlinemetanote{writes: $\text{cyclic Abelian}$}. So, this is a word we haven't really used before: cyclic Abelian.

If you meet someone and they say they're considering a cyclic Abelian group of order $n$, they mean this group, exactly this group. It's cyclic --- well, it's because the one element generates it all in a big cycle.
\end{speech}

\begin{content}[The Simplest Family: Cyclic Groups]
At the opposite structural extreme from $S_n$ lies the family of finite cyclic groups:
\begin{example}[Additive Integers Modulo $n$]\label[example]{ex:integers-mod-n}
The set of residue classes modulo $n$,
\[
\mathbb{Z}_n := \mathbb{Z} / n\mathbb{Z} = \{[0], [1], \dots, [n-1]\},
\]
equipped with addition modulo $n$ and identity element $[0]$, forms a finite Abelian group of order $n$.
\end{example}

\begin{definition}[Cyclic Group]\label[definition]{def:cyclic-group-intro}
A group $G$ is called \newterm{cyclic} if there exists an element $g \in G$ such that every element of $G$ can be written as a power (or multiple in additive notation) of $g$. In particular, $(\mathbb{Z}_n, +, [0])$ is generated by $[1]$.
\end{definition}
\end{content}

\subsection{Review of Element Orders}

\begin{speech}[00:02:04 - 00:03:20]
So, we were discussing in the second period last time... the setup of the discussion was $G$ is a group \inlinemetanote{writes: $G \text{ is group, finite}$}, a finite one, and $x$ is an element \inlinemetanote{writes: $x \in G$}. And we had proven some things already, so we had defined --- it's kind of a theorem-definition --- we had seen \inlinemetanote{writes: $\text{(1) we had seen that}$} that the order of $x$ is well-defined \inlinemetanote{writes: $\text{order of } x \text{ is well-defined}$}.

And I'll just remind you what the order is. This order of $x$ \inlinemetanote{writes: $\operatorname{Ord}(x) =$} is equal... is the minimal $m$ positive integer such that \inlinemetanote{writes: $\min_{m > 0} \text{ s.t. } x^m = e$} $x^m$ is the neutral element in $G$. So, to say it's well-defined, you have to first show that there exists some, and that it's a non-empty set and you can pick the minimal one. So, we proved all of this last time.
\end{speech}

\begin{content}[Order of a Group Element]
Let $G$ be a finite group and let $x \in G$.

\begin{definition}[Order of an Element]\label[definition]{def:order-of-element}
The \newterm{order} of $x \in G$, denoted by $\operatorname{Ord}(x)$, is the smallest positive integer $m \ge 1$ such that $x^m = e$:
\[
\operatorname{Ord}(x) := \min \{ m \in \mathbb{N}_{>0} \mid x^m = e \}.
\]
\end{definition}

\begin{steps}[Why the Order is Well-Defined in Finite Groups]
Since $G$ is finite, the infinite sequence of powers
\[
x^1, x^2, x^3, x^4, \dots
\]
cannot consist of distinct elements. By the Pigeonhole Principle, there must exist integers $j > i \ge 1$ such that $x^j = x^i$. Multiplying both sides by $x^{-i}$ yields $x^{j-i} = e$. Since $j - i > 0$, the set $\{m \in \mathbb{N}_{>0} \mid x^m = e\}$ is non-empty. By the Well-Ordering Principle of the natural numbers, a unique minimal positive integer $m$ exists.
\end{steps}
\end{content}

\subsection{The Subgroup Generated by an Element}

\begin{speech}[00:03:20 - 00:04:06]
And I think that the point where I stopped was this discussion. Another definition is I use these triangular parentheses, the angle brackets. This is by definition \inlinemetanote{writes: $\text{(2) } \langle x \rangle \stackrel{\text{def}}{=}$} the set of elements that includes the neutral element, $x$, $x^2$, all the way to $x$ to the order of $x$ minus one \inlinemetanote{writes: $\{e, x, x^2, \dots, x^{\operatorname{Ord}(x)-1}\}$}.

Okay? This is a subset of $G$ \inlinemetanote{writes: $\subset G$}. That's the definition. Sometimes in words this is the... well, let me say the words later. So, that's where we were last time. Are there questions about this? That's in my memory where we were last time.
\end{speech}

\begin{content}[The Set of Powers of an Element]
\begin{definition}[Subgroup Generated by $x$ (Preliminary Setup)]\label[definition]{def:generated-set-prelim}
For any element $x \in G$ in a finite group, we define the set of its distinct powers using angle brackets:
\[
\langle x \rangle := \bigl\{ e, x, x^2, \dots, x^{\operatorname{Ord}(x)-1} \bigr\} \subset G.
\]
\end{definition}
\end{content}

\section{Subgroups Generated by Elements}

\subsection{Verification of the Subgroup Criteria}

\begin{speech}[00:04:06 - 00:05:35]
So, the statement now, proposition, is that this subset in $G$ is a subgroup \inlinemetanote{writes: $\text{Prop. } \langle x \rangle \subset G \text{ is a subgroup}$}.

All right, and what's the proof? Well, to check it's a subgroup \inlinemetanote{writes: $\text{Proof. } \langle x \rangle \text{ is closed under } \cdot$}, I have to check several things. The first is that this subset is closed under the group product, okay? So, I have to check that.

Is this true? That means I take some person here, $x^k$, and take another one, $x^\ell$ \inlinemetanote{writes: $x^k \cdot x^l$}, and I apply the group product. And this is just... this is $x$ $k$ times, this is $x$ $\ell$ times \inlinemetanote{writes: $x \cdots x \quad x \cdots x$}, so, when I put them together, I get $x$ $k+\ell$ times \inlinemetanote{writes: $= x^{k+l}$}.

And then I get slightly worried because $k+\ell$ might be bigger than this, but of course we know that $x$ to the order \inlinemetanote{writes: $x^{\operatorname{Ord}(x)} = e$} is $e$, is the neutral element. So, I can keep removing this, so, all that matters is that's a residue modulo this number. Is that clear?

So, I can just, yeah. I hope that's clear it's closed under product, because if I come up over here, I can always keep taking... I can take out the $\operatorname{Ord}(x)$ part, because it's just the neutral element.
\end{speech}

\begin{speech}[00:05:35 - 00:06:33]
Okay, so, then we have to show that --- well, I should have said this first: $e$ is in here, but that was \inlinemetanote{writes: $e \in \langle x \rangle \checkmark$} just on the list.

And the last one we have to show is that if someone is here, then the inverse is also here \inlinemetanote{writes: $\text{if } y \in \langle x \rangle, \text{ then } y^{-1} \in \langle x \rangle$}. So, this should also be clear.

If I pick one of these elements, $x^k$, where $k$ is in this list, I can always multiply it by $x^{\operatorname{Ord}(x) - k}$ \inlinemetanote{writes: $x^k \cdot x^{\operatorname{Ord}(x)-k}$}, and it will also be in this list. Unless $k$ is zero --- there's a boundary case you can check, but it's fine also. But this times this is equal to $x^{\operatorname{Ord}(x)}$, which is $e$ \inlinemetanote{writes: $= x^{\operatorname{Ord}(x)} = e$}.

So, that means the inverse is also in this list, okay? So, all these three things are satisfied...
\end{speech}

\begin{speech}[00:06:33 - 00:07:49]
...and it is a fact that this construction, if I take $x$ and I take the neutral element and I go all the way up to the order minus one, this is a nice subset and it is a subgroup.

And moreover, if you understood this product rule where you take this and this --- and maybe I just write it: is equal to $x^r$ where $r$ is $k+\ell \bmod \text{the order}$ \inlinemetanote{writes: $= x^r \text{ where } r \equiv k+l \pmod{\operatorname{Ord}(x)}$}. If you understand this product rule, as a... maybe I write it as a corollary, corollary to the proof \inlinemetanote{writes: $\text{Cor. } \langle x \rangle \cong (\mathbb{Z}_{\operatorname{Ord}(x)}, +, 0)$}: actually, this group is isomorphic to $\mathbb{Z} / \operatorname{Ord}(x)\mathbb{Z}$ with addition and zero, okay?

And the words for this --- now I tell you the words because I didn't want to bias the issue --- is that this symbol, the words for it is the subgroup generated by the element $x$ \inlinemetanote{writes: $\langle x \rangle \subset G \text{ subgroup generated by } x$}. Okay? I didn't want to say at the beginning because I didn't want to bias the issue, but now we know it's a subgroup, I can say.
\end{speech}

\begin{content}
\begin{proposition}\label[proposition]{prop:cyclic-subgroup}
Let $G$ be a finite group and $x \in G$. Then the set $\langle x \rangle$ is a subgroup of $G$, termed the \newterm{subgroup generated by $x$}:
\[
\langle x \rangle \le G.
\]
\end{proposition}

\begin{proof}
Let $m := \operatorname{Ord}(x)$. We verify the three standard subgroup criteria:
\begin{enumerate}
    \item \textbf{Identity Element:}
    By construction, $e = x^0 \in \langle x \rangle$.

    \item \textbf{Closure under Multiplication:}
    Let $x^k, x^\ell \in \langle x \rangle$, where $0 \le k, \ell < m$. The product in $G$ satisfies:
    \[
    x^k \cdot x^\ell = \underbrace{(x \cdot \dots \cdot x)}_{k \text{ times}} \cdot \underbrace{(x \cdot \dots \cdot x)}_{\ell \text{ times}} = x^{k+\ell}.
    \]
    By Euclidean division, we can write $k + \ell = q \cdot m + r$ with $0 \le r < m$, so $r \equiv k + \ell \pmod m$. Since $x^m = e$, we compute:
    \[
    x^{k+\ell} = (x^m)^q \cdot x^r = e^q \cdot x^r = x^r.
    \]
    Since $0 \le r < m$, we have $x^r \in \langle x \rangle$.

    \item \textbf{Closure under Inverses:}
    Let $y \in \langle x \rangle$, so $y = x^k$ for some $0 \le k < m$.
    \begin{itemize}
        \item \textbf{Case $k = 0$:} then $y = e$ and $y^{-1} = e \in \langle x \rangle$.
        \item \textbf{Case $1 \le k < m$:} consider $y' = x^{m - k}$. Since $0 < m - k < m$, we have $y' \in \langle x \rangle$. Multiplying yields:
        \[
        y \cdot y' = x^k \cdot x^{m-k} = x^m = e.
        \]
        Thus, $y^{-1} = x^{m-k} \in \langle x \rangle$.
    \end{itemize}
\end{enumerate}
Hence, $\langle x \rangle$ is a subgroup of $G$.
\end{proof}

\begin{corollary}\label[corollary]{cor:cyclic-subgroup-isomorphism}
Let $G$ be a finite group and $x \in G$ with $m = \operatorname{Ord}(x)$. Then the subgroup $\langle x \rangle$ is isomorphic to the additive cyclic group $(\mathbb{Z}_m, +, [0])$:
\[
\langle x \rangle \cong (\mathbb{Z}_m, +, [0]).
\]
In particular, $|\langle x \rangle| = \operatorname{Ord}(x)$.
\end{corollary}

\begin{steps}[Why the Map Is an Isomorphism]
The assignment $\psi \colon \mathbb{Z}_m \to \langle x \rangle$ defined by $\psi([k]) := x^k$ is well-defined because $x^m = e$. It is surjective by definition of $\langle x \rangle$, and since both sets have cardinality $m$, it is bijective. Furthermore,
\[
\psi([k] + [\ell]) = \psi([k + \ell]) = x^{k+\ell} = x^k \cdot x^\ell = \psi([k]) \cdot \psi([\ell]),
\]
which proves that $\psi$ is a group isomorphism.
\end{steps}
\end{content}

\section{Lagrange's Theorem}

\subsection{Divisibility of Orders in Symmetric Groups}

\begin{speech}[00:07:49 - 00:09:18]
All right, so, like in our example for $S_3$, we had some elements like these two-cycles, this element \inlinemetanote{writes: $S_3$, $(1\,2)(3) \implies \operatorname{ord} = 2$}, and I said sometimes I write this, sometimes I don't, it's just, you know, however I feel. But often people don't write the singletons.

If I take this element, we saw it has order $2$. So, this has order $2$, and this \inlinemetanote{writes: $(1\,2\,3) \implies \operatorname{ord} = 3$} has order $3$.

And one of the things you notice here is: what is the order of this group? It's $3$ factorial, which is $6$ \inlinemetanote{writes: $|S_3| = 3! = 6$}. So, these orders, we find here the orders of the elements, at least in this example, orders of elements seem to divide the order of the group \inlinemetanote{writes: $\text{Orders of elem. } \mid |S_3|$}.

At least in this case. So, you can wonder whether this is always the case. So, this is kind of an interesting question. As I said, like this kind of question is a question about finite groups. Because if $G$ has infinite order, we can't really talk about... well, if $G$ was a group with infinite order, the elements may not have finite order, and in any case we can't talk about this division property. That's why I say we're focusing on finite groups.
\end{speech}

\begin{content}[Motivating Example: Element Orders in \texorpdfstring{$S_3$}{S3}]
Consider the symmetric group on three letters, $S_3$, of order:
\[
|S_3| = 3! = 6.
\]
Evaluating the orders of representative elements:
\begin{align*}
\sigma = (1\ 2)(3) &\implies \operatorname{Ord}(\sigma) = 2, && \text{and } 2 \mid 6, \\
\tau = (1\ 2\ 3) &\implies \operatorname{Ord}(\tau) = 3, && \text{and } 3 \mid 6, \\
\operatorname{id} &\implies \operatorname{Ord}(\operatorname{id}) = 1, && \text{and } 1 \mid 6.
\end{align*}
In every case, the order of the element divides the order of the group:
\[
\operatorname{Ord}(x) \bigm| |S_3| \quad \forall x \in S_3.
\]
\end{content}

\subsection{Statement of Lagrange's Theorem}

\begin{speech}[00:09:18 - 00:10:50]
So, anyway, in this case we find that the order of the elements divide the order of the group, so, one can wonder: is this always the case? \inlinemetanote{writes: $\text{Is this always true?}$}

Specifically, if $x$ is an element of $G$, a finite group, does the order of $x$ divide the order of $G$? \inlinemetanote{writes: $\text{(If } x \in G \text{ finite group, does } \operatorname{Ord}(x) \mid |G|\text{?)}$} That's a very specific question, and the answer is going to be yes \inlinemetanote{writes: $\text{YES}$}.

And actually the proof is that a much bigger property, a much more general property holds, and that's the main topic of today's lecture, which is really in some sense usually the first theorem in finite group theory: it's a theorem of Joseph-Louis Lagrange \inlinemetanote{writes: $\text{J. L. Lagrange}$}.

And it's that if $H$ is a subgroup of $G$ --- $H$ is a subgroup and $G$ is finite \inlinemetanote{writes: $\text{Thm: If } H \subset G \text{ is a subgroup and } G \text{ is finite}$} --- so that's the hypothesis: we have a finite group and a subgroup. I don't have to say the subgroup is finite, because once the group is finite, the subgroup is also finite. So, then, this is what Joseph-Louis says: the order of $H$ always divides the order of $G$ \inlinemetanote{writes: $\text{then } |H| \mid |G|$}. Okay?
\end{speech}

\begin{content}
\begin{question}[Element Order Divisibility]\label[question]{q:order-divides-group-order}
If $G$ is a finite group and $x \in G$, does $\operatorname{Ord}(x)$ always divide $|G|$?
\end{question}

\begin{answer}[to \cref{q:order-divides-group-order}]\label[answer]{ans:order-divides-group-order}
\textcolor{ForestGreen}{\textbf{Yes}}. In fact, this follows as an immediate corollary of a far more general theorem governing all subgroups of a finite group.
\end{answer}

\begin{theorem}[Lagrange's Theorem]\label[theorem]{thm:lagrange}
Let $G$ be a finite group and let $H \le G$ be a subgroup of $G$. Then the order of $H$ divides the order of $G$:
\[
|H| \bigm| |G|.
\]
\end{theorem}

\begin{steps}[Connection to Element Orders]
By \cref{cor:cyclic-subgroup-isomorphism}, for any $x \in G$, the generated subgroup $H = \langle x \rangle$ has order $|H| = \operatorname{Ord}(x)$. Therefore, \cref{thm:lagrange} immediately implies:
\[
\operatorname{Ord}(x) \bigm| |G| \quad \forall x \in G.
\]
\end{steps}
\end{content}

\subsection{Historical Context of the Theorem}

\begin{speech}[00:10:50 - 00:12:05]
This is usually the very first actual theorem about finite groups. And I will prove it; the proof is simple, it's conceptual.

One interesting thing is that he lived in the 18th century, and the definition of a group came like half a century after he passed away, so you can wonder: how did he prove this theorem? And, you know, of course he didn't prove this theorem. He proved something about symmetric groups that people have interpreted to be this theorem.

But this is the same Lagrange that's responsible for Lagrangian mechanics and the calculus of variations, which is a fantastic topic. Whether that's covered in the math curriculum, I don't know. Part of my life was as physicists and I like that part. That's a great topic. So, he's done a lot of amazing stuff.

But you've taken this stuff like Lagrangian mechanics at some point? Calculus of variations? Do you take that? Yeah. Those are really great subjects.
\end{speech}

\begin{aiNote}[Beyond the Lecture: Historical Context of Lagrange's Theorem]
% [PERSISTENT, ADDITION]
Joseph-Louis Lagrange (1736--1813) published his foundational work on the solvability of algebraic equations, \emph{R{\'e}flexions sur la r{\'e}solution alg{\'e}brique des {\'e}quations}, in 1770. In that treatise, Lagrange analyzed permutations of roots of polynomial equations and showed that the number of different values that a rational function of $n$ variables can take under all $n!$ permutations must divide $n!$. The abstract definition of a group was not formulated until the work of Arthur Cayley in 1854, more than four decades after Lagrange's death. Modern algebraists later abstracted Lagrange's permutation calculation into the universal statement $|H| \mid |G|$ for any subgroup of a finite group.
\end{aiNote}

\section{Proof of Lagrange's Theorem}

\subsection{The Subgroup Equivalence Relation}

\begin{speech}[00:12:05 - 00:13:17]
Okay, so, we're going to prove this. Proof \inlinemetanote{writes: $\text{Proof.}$}. Is there any question about this theorem?

And you see this theorem certainly implies this statement, because if I have an element, the order of this element is the size of that subgroup. So, in particular, if every subgroup divides the order of $G$, the order always divides the order of $G$, okay? But then this is a much more general result.

So, how are we going to prove it? It's an extremely nice proof, and yeah. So, the proof is, we start by defining an equivalence relation on the set $G$ \inlinemetanote{writes: $\text{We start by defining an equivalence relation on set } G.$}. And then we will count... we will count $G$ by counting up the size of the equivalence classes.

And this equivalence relation is just incredibly important for group theory. So, we will define it.
\end{speech}

\begin{speech}[00:13:17 - 00:14:52]
The notion of a group, while I erase I say, is that, you know, it has many fewer axioms than a ring, and in some sense it's a more primitive notion in mathematics, which means that it's more widely applicable and in some sense more fundamental than the notion of a ring. The ring is motivated by number systems, so somehow it is a natural choice for the first topic in the course, since most people have familiarity with number systems. But the group is a more ubiquitous notion, because of course it has fewer axioms, so you can find it more generally. And it's motivated --- I mean, it has many motivations, but a very strong motivation about the group is about symmetries of geometric objects. It's like a totally different motivation than a ring, which is about number systems that you encounter in school.

So, while they're in the same class, and in some sense one thinks about them algebraically in the same space, the motivations that led to them are pretty much different, that I would say. But they're both really amazingly interesting, intricate topics. And a lot of the algebraic moves are similar. Maybe that justifies having them in the same course. Sometimes the course is taught with groups first and then rings. I've taught it with rings first, then groups. It doesn't really matter, of course.
\end{speech}

\begin{note}[Board Transition and Conceptual Motivations]
The lecturer erases the board and reflects on the broader architecture of abstract algebra: while rings capture arithmetic properties modeled after number systems ($\mathbb{Z}, \mathbb{Q}, \mathbb{R}$), groups possess fewer axioms, making them more fundamental and universally applicable to transformations and geometric symmetries.
\end{note}

\begin{speech}[00:14:52 - 00:16:00]
All right, so, for the proof of this theorem, I'm going to define an equivalence relation. So, that means if you have $x$ and $y$ in $G$ \inlinemetanote{writes: $x, y \in G$}, I have to tell you how to define this relation. The definition \inlinemetanote{writes: $x \sim y \stackrel{\text{def}}{\iff}$} is that we say $x$ is related to $y$ if and only if --- the definition --- if there exists an $h$ in the subgroup \inlinemetanote{writes: $\exists h \in H$} such that $xh = y$ \inlinemetanote{writes: $\text{s.t. } xh = y$}. All right?

Now, you could say: why do I put $h$ on the right side? Could I put it on the left side? And you could do that. That'd be your own equivalence relation, and it's totally equivalent for this proof, and the proof will be fine. But since the groups are not commutative, not necessarily commutative, I have to decide, and I've decided to put it on the right, okay? That's the first comment.
\end{speech}

\begin{content}[The Equivalence Relation Induced by \texorpdfstring{$H$}{H}]
Let $G$ be a group and $H \le G$ a subgroup. For elements $x, y \in G$, we define the relation $\sim$ by:
\begin{equation}\label{eq:coset-equiv-relation}
x \sim y \stackrel{\text{def}}{\iff} \exists h \in H \text{ such that } xh = y.
\end{equation}

\begin{steps}[Left vs. Right Coset Conventions]
Because $G$ is not assumed to be Abelian, the equation $xh = y$ places $h$ on the right, defining what is conventionally known as the \newterm{left coset} relation (since $y \in xH$). Placing $h$ on the left ($hx = y$) yields an entirely symmetric theory of right cosets ($Hx$). The proof of Lagrange's Theorem works identically under either convention.
\end{steps}
\end{content}

\subsection{Verifying the Equivalence Relation Axioms}

\begin{speech}[00:16:00 - 00:17:28]
So, then, I have to check for an equivalence relation, I have to check the properties. And normally I leave this to you, but here this is so fundamental I'm going to do it. It's probably a mistake, but nevertheless.

The first property is that $x$ is related to itself \inlinemetanote{writes: $\bullet\ x \sim x$}. And that's if and only if I can find an $h$ in this group that satisfies that \inlinemetanote{writes: $\iff \exists h \in H \text{ s.t. } xh = x$}, and which one is that? Yeah, neutral element, because $e$ has to be here, because it's a subgroup, and then of course I have this equation \inlinemetanote{writes: $e \in H \quad x e = x \checkmark$}, so that's fine.

Now, that's the reflexive property. Then the symmetric property is if $x$ is related to $y$, then $y$ is related to $x$ \inlinemetanote{writes: $\bullet\ \text{If } x \sim y \implies y \sim x$}. So, we have to prove this. $x$ is related to $y$ means that there exists an $h$ such that $xh = y$ \inlinemetanote{writes: $x \sim y \implies \exists h \in H \text{ s.t. } xh = y$}. And then what do I write on the then-side? If I have this, I can just multiply by $h$ inverse, right? So, then, I get... sorry... so, I can multiply by $h$ inverse, which will make this $h$ go away and put $h$ inverse there \inlinemetanote{writes: $\text{then } x = yh^{-1}$}. And then this means that $y$ is related to $x$. I can write this again if you really want to fit the definition \inlinemetanote{writes: $yh^{-1} = x$}: this says that $y$ is related to $x$ because if you have an element of $H$, the inverse is also in $H$ \inlinemetanote{writes: $h^{-1} \in H \implies y \sim x \checkmark$}, that's a property of being in a subgroup, right?
\end{speech}

\begin{speech}[00:17:28 - 00:19:02]
And the third one is transitivity, which says that if I have $x$ related to $y$ and $y$ related to $z$, then I have to show that $x$ is related to $z$ \inlinemetanote{writes: $\bullet\ \text{If } x \sim y \text{ and } y \sim z \implies x \sim z$}. And if I do this using the definition, this means that there exists some $h_1$ that satisfies this equation \inlinemetanote{writes: $xh_1 = y$}, and this one means there's some $h_2$ that satisfies this equation \inlinemetanote{writes: $yh_2 = z$}. And then I have to find this relation equation, and you can imagine what I do: I just take here $x h_1 h_2 = z$ \inlinemetanote{writes: $\implies x(h_1 h_2) = z$}, all right?

And why does this work? It's because it's a subgroup, so, $H$ is closed under product \inlinemetanote{writes: $h_1 h_2 \in H \text{ as } H \text{ is closed under product } \checkmark$}. So, that's the proof that it's an equivalence relation.

And you see the proof is perfect in a way that I need to check three properties to be an equivalence relation, and those were exactly the same three properties that $H$ is a subgroup. I've used everything exactly: the existence of the neutral element is there, it's closed under inverse, and it's closed under product. So, in some sense, the statement that this definition is... on one hand, on this side it's an equivalence relation, and on that side $H$ is a subgroup, they're just perfectly aligned. All right, so, we have proven that.
\end{speech}

\begin{content}
\begin{proposition}\label[proposition]{prop:equiv-relation-subgroup}
The relation $\sim$ defined in \eqref{eq:coset-equiv-relation} is an equivalence relation on $G$.
\end{proposition}

\begin{proof}
We verify the three defining axioms of an equivalence relation:
\begin{itemize}
    \item \textbf{Reflexivity ($x \sim x$):}
    We require an element $h \in H$ such that $xh = x$. Since $H$ is a subgroup, the identity element satisfies $e \in H$, and $xe = x$. Thus, $x \sim x$.

    \item \textbf{Symmetry ($x \sim y \implies y \sim x$):}
    Suppose $x \sim y$. By definition, there exists $h \in H$ such that $xh = y$. Multiplying on the right by $h^{-1} \in G$ yields:
    \[
    x = yh^{-1} \iff yh^{-1} = x.
    \]
    Since $H$ is a subgroup, $h \in H \implies h^{-1} \in H$. Therefore, setting $h' := h^{-1} \in H$ gives $yh' = x$, which proves $y \sim x$.

    \item \textbf{Transitivity ($x \sim y \text{ and } y \sim z \implies x \sim z$):}
    Suppose $x \sim y$ and $y \sim z$. Then there exist elements $h_1, h_2 \in H$ such that:
    \[
    xh_1 = y \quad \text{and} \quad yh_2 = z.
    \]
    Substituting the expression for $y$ into the second equation yields:
    \[
    (xh_1)h_2 = z \implies x(h_1 h_2) = z.
    \]
    Since $H$ is closed under multiplication, $h_3 := h_1 h_2 \in H$. Thus, $x \sim z$.
\end{itemize}
Hence, $\sim$ is an equivalence relation on $G$.
\end{proof}

\begin{steps}[Correspondence between Subgroup Axioms and Equivalence Axioms]
Notice the exact structural alignment highlighted by the lecturer:
\begin{align*}
e \in H &\iff \text{Reflexivity of } \sim \\
h \in H \implies h^{-1} \in H &\iff \text{Symmetry of } \sim \\
h_1, h_2 \in H \implies h_1 h_2 \in H &\iff \text{Transitivity of } \sim
\end{align*}
The definition of the relation translates the three subgroup axioms directly into the three equivalence relation axioms.
\end{steps}
\end{content}

\subsection{Partitioning into Equivalence Classes and the Counting Formula}

\begin{speech}[00:19:02 - 00:20:33]
Okay, so, now we have a picture. Here is the group $G$, the finite group, so there's some people here \inlinemetanote{draws a contour for $G$ with dots inside}. And I have an equivalence relation. Whenever you have an equivalence relation, it divides the set into equivalence classes.

So, our equivalence relation divides it into equivalence classes. Maybe here's one equivalence class, $E_1$ \inlinemetanote{draws boundary for $E_1$}, and then there's some equivalence class here, $E_2$ \inlinemetanote{draws boundary for $E_2$}, and then there's another equivalence class, $E_3$ \inlinemetanote{draws boundary for $E_3$}, I don't know. There's finitely many --- it's a finite set, so it's finitely many equivalence classes.

So, we can write an equation, and this equation you can always write for equivalence classes on finite sets, and it's an equation that is so naive you think it couldn't possibly be useful: the equation is that the size of $G$ is equal to the sum from $i=1$ to $k$ \inlinemetanote{writes: $|G| = \sum_{i=1}^k |E_i|$}, where $k$ is the number of equivalence classes \inlinemetanote{writes: $k = \text{\# of equiv classes}$} --- there are some finite number of equivalence classes, it's a finite group --- you sum over the number of equivalence classes and it's the size of the equivalence class.

So, that's the equation. It's a pretty naive equation, but as I said, it looks so silly it seems hard to believe that's useful, but it's just saying you take this $G$, equivalence relation makes some equivalence classes, and you account for everybody in $G$, and that equation just says you've accounted for everybody exactly once.
\end{speech}

\begin{content}[Partition of \texorpdfstring{$G$}{G} into Equivalence Classes]
Since $\sim$ is an equivalence relation, the fundamental theorem of equivalence relations ensures that $G$ is partitioned into disjoint, non-empty equivalence classes:
\[
G = E_1 \;\dot{\cup}\; E_2 \;\dot{\cup}\; \dots \;\dot{\cup}\; E_k,
\]
where $k$ is the total number of distinct equivalence classes (cosets).

\definecolor{figBlue}{HTML}{076090} % RGB (7, 96, 144): petrol blue (MidnightBlue)
\definecolor{figGreen}{HTML}{3C8252} % RGB (60, 130, 82): dark green (ForestGreen!70!black)
\definecolor{figOrange}{HTML}{CC7F2B} % RGB (204, 127, 43): ochre orange (BurntOrange!80!black)
\definecolor{figGray}{HTML}{666666} % RGB (102, 102, 102): dark gray
\definecolor{figLightGray}{HTML}{A6A6A5} % RGB (166, 166, 165): light gray
\begin{hidden}
- Canvas: no scale, cm; outer contour for G centered around (0,0), spans x in [-4.0, 4.5], y in [-2.5, 2.5]
- Order: (1) Fill and outline of G; (2) Three partitioned cells E_1, E_2, E_3 with distinct outlines and light fills;
  (3) Dots representing elements; (4) Labels E_1, E_2, E_3, G, dots \dots
- Geometry:
  G outer domain: rounded rectangle / smooth path x in [-3.8, 3.8], y in [-2.2, 2.2]
  Cell E_1 (top left): ellipse centered at (-1.8, 0.9), radii 1.3 and 0.8, rotated 15 deg
  Cell E_2 (bottom left): ellipse centered at (-1.5, -0.9), radii 1.4 and 0.8, rotated -10 deg
  Cell E_3 (right): ellipse centered at (1.5, 0.2), radii 1.6 and 1.2, rotated -5 deg
- Labels:
  "G" placed bottom right at (3.2, -1.6)
  "E_1" centered in E_1 at (-1.8, 1.2)
  "E_2" centered in E_2 at (-1.5, -0.6)
  "E_3" centered in E_3 at (1.5, 0.6)
- Collisions: Ensure labels don't collide with dots; text drawn last.
\end{hidden}
\begin{center}
\begin{tikzpicture}[>=Stealth]
    % (1) Outer Boundary of Group G
    \draw[thick, figGray, fill=figLightGray!10] 
        (-3.8, -1.8) .. controls (-4.0, 0.5) and (-2.5, 2.2) .. (0, 2.2)
                     .. controls (2.5, 2.2) and (4.0, 1.0) .. (4.0, -0.5)
                     .. controls (4.0, -2.0) and (2.0, -2.2) .. (0, -2.2)
                     .. controls (-2.0, -2.2) and (-3.5, -2.2) .. cycle;

    % (2) Equivalence Classes E_1, E_2, E_3 (Cosets)
    \draw[thick, figBlue, fill=figBlue!8] (-1.8, 0.8) ellipse (1.3cm and 0.8cm);
    \draw[thick, figGreen, fill=figGreen!8] (-1.5, -1.0) ellipse (1.4cm and 0.75cm);
    \draw[thick, figOrange, fill=figOrange!8] (1.6, 0.1) ellipse (1.6cm and 1.3cm);

    % (3) Discrete elements (dots) inside classes
    \fill (-2.2, 0.9) circle (1.5pt) (-1.6, 0.6) circle (1.5pt) (-1.4, 1.0) circle (1.5pt);
    \fill (-1.9, -1.1) circle (1.5pt) (-1.4, -0.8) circle (1.5pt) (-1.0, -1.2) circle (1.5pt);
    \fill (0.9, 0.2) circle (1.5pt) (1.4, -0.3) circle (1.5pt) (1.8, 0.5) circle (1.5pt) (2.4, -0.1) circle (1.5pt);

    % Other scattered dots in G
    \fill (3.0, 1.2) circle (1.5pt) (2.8, -1.2) circle (1.5pt) (-3.0, -0.2) circle (1.5pt);

    % (4) Text last
    \node[anchor=north east, font=\large, figGray] at (3.6, -1.6) {$G$};
    \node[figBlue, font=\bfseries] at (-1.8, 1.2) {$E_1$};
    \node[figGreen, font=\bfseries] at (-1.5, -0.6) {$E_2$};
    \node[figOrange, font=\bfseries] at (1.6, 0.9) {$E_3$};
    \node[figGray, font=\Large] at (0.0, -0.2) {$\dots$};
\end{tikzpicture}
\end{center}

Since the equivalence classes are pairwise disjoint and their union exhausts $G$, the sum of their individual sizes accounts for every element in $G$ exactly once:
\begin{equation}\label{eq:counting-classes}
|G| = \sum_{i=1}^k |E_i|,
\end{equation}
where $k$ denotes the total number of equivalence classes.
\end{content}

\subsection{Deducing Lagrange's Theorem from the Cardinality Claim}

\begin{speech}[00:20:33 - 00:21:22]
This fellow Joseph-Louis, he was actually born Italian, and his name at birth was Giuseppe Luigi, as far as I can tell. But he became French somehow in his life. His name was always Lagrange as far... but he was born in Turin.

The standard place in the undergraduate curriculum where you learn calculus of variations would be a mathematical class on classical mechanics, which is often taught probably in the physics department, I don't know. But that's probably his greatest work, Lagrange's greatest work. This is not his greatest work. But it's very well known.
\end{speech}

\begin{speech}[00:21:22 - 00:22:28]
All right, so, what are we going to do here? We have to prove this, and so, the observation, which is an incredibly simple observation... So, the claim, which will prove everything, is that for all equivalence classes, the size of this equivalence class is actually the size of $H$ \inlinemetanote{writes: $\boxed{\text{Claim: } \forall i, \ |E_i| = |H|}$}. All right?

So, that sum is actually just adding up the order of $H$. And so, of course, that proves it. Claim: if I prove this --- so, I have to prove this claim --- then Lagrange's Theorem follows \inlinemetanote{writes: $\text{Then Lagrange's Theorem follows}$}, since we rewrite that equation as: the size of $G$ is equal to the sum $i=1$ to $k$ of $|E_i|$ \inlinemetanote{writes: $|G| = \sum_{i=1}^k |E_i|$}, but all of these are just the size of $H$, so this is actually $k \times |H|$ \inlinemetanote{writes: $= k \cdot |H|$}.

So, that shows that, as a consequence, the order of $H$ divides the order of $G$ \inlinemetanote{writes: $\implies |H| \mid |G| \ \ \square$}. So, that's the end of the theorem, and I just have to prove this. Okay?
\end{speech}

\begin{content}
\begin{claim}[Cardinality of Equivalence Classes]\label[claim]{clm:class-cardinality}
Every equivalence class $E_i$ has the same size as the subgroup $H$:
\[
\boxed{\forall i \in \{1, \dots, k\} \colon \ |E_i| = |H|.}
\]
\end{claim}

\begin{proof}[Proof of \cref{thm:lagrange} assuming \cref{clm:class-cardinality}]
Assuming that \cref{clm:class-cardinality} holds, we substitute $|E_i| = |H|$ into the counting equation \eqref{eq:counting-classes}:
\begin{align}
|G| &= \sum_{i=1}^k |E_i| \nonumber \\
    &= \sum_{i=1}^k |H| \nonumber \\
    &= k \cdot |H|. \label{eq:lagrange-index-formula}
\end{align}
Since $k$ is an integer (the number of cosets), equation \eqref{eq:lagrange-index-formula} directly implies:
\[
|H| \bigm| |G|.
\]
This completes the proof of Lagrange's Theorem, modulo proving \cref{clm:class-cardinality}.
\end{proof}

\begin{steps}[The Concept of Group Index]
The integer $k$, representing the number of distinct cosets of $H$ in $G$, is called the \newterm{index} of $H$ in $G$, conventionally written as:
\[
[G : H] := k = \frac{|G|}{|H|}.
\]
\end{steps}
\end{content}

\subsection{Proof of the Cardinality Claim}

\begin{speech}[00:22:28 - 00:23:53]
So, let us see: how are we going to prove that? Let $E = E_i$ be an equivalence class in $G$ \inlinemetanote{writes: $\text{Let } E = E_i \text{ be an equiv. class in } G$}. So, we have some equivalence class in $G$. Here's the picture of $G$. Here is $E$ \inlinemetanote{draws set $G$ and subset $E$}, and I need to show that the cardinality of $E$ is exactly $|H|$.

The first thing I do is I reach into $E$ and I pick a person in it, which of course I call $x$, right? \inlinemetanote{draws point $x$ in $E$ and writes: $x \in E$} So, I have $x$ in $E$. I can do that, because it's an equivalence class, it's not empty; we don't consider the empty equivalence class.

Okay, so, then I would like to make a map from $H$ to $E$ \inlinemetanote{writes: $\phi \colon H \to E$}, and I define this: $\phi(h) = \dots$ \inlinemetanote{writes: $\phi(h) =$} And what do you think I'm going to write here?

I have to write something that's definitely in $E$, and all I know about $E$, it has $x$ in it. So, you have to remember what this equivalence relation... So, yeah, what would you propose I write? I can't... yeah.
\end{speech}

\begin{student}[Student Answer]
$xh$
\end{student}

\begin{speech}[00:23:54 - 00:25:12]
$xh$ \inlinemetanote{writes: $xh$}. I have to write something.

And whatever I write has to be in $E$. And the things that are in $E$, remember: $x$ is equivalent to $y$ if and only if we have this equation, $xh = y$ for some $h$ in $H$ \inlinemetanote{writes: $x \sim y \iff \exists h \in H \text{ s.t. } xh = y$}. So, these things for sure are equivalent to $x$.

Is that clear? This is more or less the only thing you could possibly write in this situation.

So, that gives a map, so that at least $\phi$ maps $H$ to $E$ \inlinemetanote{writes: $\text{At least } \phi \colon H \to E$}. That means if you put in any $h$ in here, you get something in the equivalence class.

So, now we want to prove --- the last step is we should prove, we will prove, let's be confident here --- prove that $\phi$ is a bijection \inlinemetanote{writes: $\text{We will prove that } \phi \text{ is a bijection}$}.
\end{speech}

\begin{content}[Construction of the Bijection between \texorpdfstring{$H$}{H} and \texorpdfstring{$E$}{E}]
Let $E = E_i$ be an arbitrary equivalence class of $\sim$. Since equivalence classes are non-empty, we choose an arbitrary element $x \in E$.

We define the mapping:
\begin{equation}\label{eq:phi-map-def}
\phi \colon H \to E, \quad \phi(h) := xh.
\end{equation}

\definecolor{figBlue}{HTML}{076090} % RGB (7, 96, 144): petrol blue (MidnightBlue)
\definecolor{figGray}{HTML}{666666} % RGB (102, 102, 102): dark gray
\definecolor{figLightGray}{HTML}{A6A6A5} % RGB (166, 166, 165): light gray (gray!70)
\begin{hidden}
- Canvas: no scale, cm; outer domain G and inner subset E centered at (0, 0); x in [-3.5, 3.5], y in [-2.0, 2.0]
- Order: (1) Domain G boundary and fill; (2) Subset E boundary and fill; (3) Chosen element x; (4) Labels G, E, x
- Coordinates:
  Domain G: smooth contour path around (0, 0), width ~6.2cm, height ~3.2cm
  Subset E: ellipse centered at (0.2, 0.1), radii 1.4cm and 0.9cm, fill figBlue!10, draw figBlue
  Point x: fill at (0.2, 0.1)
- Labels (normal size, 0.20 per character, height 0.45):
  "G" anchor north east at (3.0, -1.1): x in [2.8, 3.0], y in [-1.55, -1.1]
  "E" centered at (-0.6, 0.5): x in [-0.7, -0.5], y in [0.28, 0.72]
  "x" anchor west at (0.35, 0.1): x in [0.35, 0.55], y in [-0.12, 0.32]
- Collisions: x is centered inside E, 1.05cm from E's boundary; E is well inside G; text drawn last.
\end{hidden}
\begin{center}
\begin{tikzpicture}[>=Stealth]
    % (1) Outer boundary G
    \draw[thick, figGray, fill=figLightGray!10]
        (-3.0, -1.4) .. controls (-3.3, 0.2) and (-2.0, 1.6) .. (0, 1.6)
                     .. controls (2.0, 1.6) and (3.3, 0.8) .. (3.2, -0.4)
                     .. controls (3.0, -1.5) and (1.5, -1.6) .. (0, -1.6)
                     .. controls (-1.5, -1.6) and (-2.7, -1.6) .. cycle;

    % (2) Equivalence class E
    \draw[thick, figBlue, fill=figBlue!10] (0.2, 0.1) ellipse (1.4cm and 0.9cm);

    % (3) Chosen representative x
    \fill (0.2, 0.1) circle (2pt);

    % (4) Text last
    \node[anchor=north east, font=\large, figGray] at (3.0, -1.1) {$G$};
    \node[figBlue, font=\bfseries] at (-0.6, 0.5) {$E$};
    \node[anchor=west] at (0.35, 0.1) {$x$};
\end{tikzpicture}
\end{center}

\begin{steps}[Well-Definedness of $\phi$]
For every $h \in H$, setting $y := xh$ gives $xh = y$ with $h \in H$, which satisfies the definition of the equivalence relation $x \sim y$. Since $x \in E$, it follows that $xh \in E$. Thus, $\phi(H) \subseteq E$, confirming that $\phi \colon H \to E$ is well-defined.
\end{steps}
\end{content}

\begin{speech}[00:25:12 - 00:25:57]
So, if we do that, then we're happy, because we need to show that the cardinality of $E$ is the same as the cardinality of $H$, and if I show this map is a bijection, then I will have shown that.

All right, so, for this we need to start... we need to prove that $\phi$ is surjective \inlinemetanote{writes: $\bullet\ \phi \text{ is surjective}$}. And this is true because anybody in $E$ is someone that's related to $x$, and all people related to $x$ are of this form. So, it's surjective by definition. And this is true by definition of this equivalence relation \inlinemetanote{writes: $\text{(True by def of } \sim)$}, meaning the only way to be in this equivalence class of $x$ is to satisfy such a relation for some $h$.
\end{speech}

\begin{speech}[00:25:57 - 00:27:37]
Okay, then we have to prove $\phi$ is injective \inlinemetanote{writes: $\bullet\ \phi \text{ is injective}$}, and then we will be done.

Okay, so, we have to prove this. Why? Why would this be true? Injectivity means you have to study the equation $\phi(h_1) = \phi(h_2)$ \inlinemetanote{writes: $\phi(h_1) = \phi(h_2)$}. We have to study that equation, right? So, let's see: suppose $\phi(h_1) = \phi(h_2)$, where $h_1$ and $h_2$ are in $H$ \inlinemetanote{writes: $h_1, h_2 \in H$}.

We suppose this is true. Well, what do we do next? We write down the definitions. This thing is $xh_1$, and that should equal $xh_2$ \inlinemetanote{writes: $\iff xh_1 = xh_2$}. All right. But this is an equation that holds in the group.

So, we now use the cancellation property. By the cancellation property, we can get rid of this \inlinemetanote{cancels $x$ from both sides}. This implies $h_1 = h_2$ \inlinemetanote{writes: $\iff h_1 = h_2$}, and that's injectivity.

So, we finished the proof. Yeah, that's the whole proof. It's an extremely simple proof, and if you understand it, you'll never forget it.

When you first see this, it's a very strange property that if you have these finite groups and you assemble these subgroups, you can try as hard as you want, but you'll never be able to make a subgroup that doesn't have order which divides the group.

In particular, if your group has prime order, it's extremely hard to make subgroups, because there's very few numbers that divide prime numbers. So, this is going to be used over and over again in the study of finite groups.
\end{speech}

\begin{content}[Verification of Bijectivity and Completion of the Proof]
\begin{proof}[Proof of \cref{clm:class-cardinality}]
To prove $|E| = |H|$, we show that the map $\phi \colon H \to E$, $\phi(h) := xh$, is a bijection:
\begin{enumerate}
    \item \textbf{Surjectivity:}
    Let $y \in E$. Since $E$ is the equivalence class of $x$, we have $x \sim y$. By definition of $\sim$, there exists an element $h \in H$ such that:
    \[
    xh = y \iff \phi(h) = y.
    \]
    Hence, every element $y \in E$ has a pre-image in $H$, proving that $\phi$ is surjective.

    \item \textbf{Injectivity:}
    Let $h_1, h_2 \in H$ satisfy $\phi(h_1) = \phi(h_2)$. By definition of $\phi$:
    \[
    xh_1 = xh_2.
    \]
    Multiplying on the left by $x^{-1} \in G$ (left cancellation in the group $G$) gives:
    \[
    x^{-1}(xh_1) = x^{-1}(xh_2) \implies (x^{-1}x)h_1 = (x^{-1}x)h_2 \implies h_1 = h_2.
    \]
    Thus, $\phi$ is injective.
\end{enumerate}
Since $\phi$ is both injective and surjective, it is a bijection. Therefore, $|E| = |H|$. This completes the proof of \cref{clm:class-cardinality}, and consequently proves Lagrange's Theorem (\cref{thm:lagrange}).
\end{proof}

\begin{corollary}[Subgroups of Groups of Prime Order]\label[corollary]{cor:prime-order-subgroups}
Let $G$ be a finite group whose order $|G| = p$ is a prime number. Then the only subgroups of $G$ are the trivial subgroup $\{e\}$ and $G$ itself.
\end{corollary}

\begin{steps}[Classification of Groups of Prime Order]
By Lagrange's Theorem, if $H \le G$, then $|H| \mid p$. Since $p$ is prime, the only positive divisors of $p$ are $1$ and $p$. Thus:
\begin{itemize}
    \item \textbf{Case $|H| = 1$:} then $H = \{e\}$.
    \item \textbf{Case $|H| = p$:} then $H = G$.
\end{itemize}
Furthermore, choosing any non-identity element $x \in G \setminus \{e\}$, its generated subgroup $\langle x \rangle$ must have order $> 1$, forcing $|\langle x \rangle| = p$, and thus $\langle x \rangle = G$. Consequently, every group of prime order is cyclic: $G \cong \mathbb{Z}_p$.
\end{steps}
\end{content}

\begin{note}[Board Transition]
The lecturer erases the middle and lower blackboards to transition from the theoretical proof of Lagrange's Theorem to concrete families of groups: symmetric groups, cycle notation, and rotational symmetry groups of geometric solids.
\end{note}

\section{Symmetric Groups and Geometric Symmetries}

\subsection{Cycle Decomposition in \texorpdfstring{$S_n$}{Sn} and Integer Partitions}

\begin{speech}[00:28:02 - 00:29:27]
Okay, I wanted to do some examples of groups, because that's the fun part of --- at least for me, the fun part of group theory is this rich source of examples.

Okay, so, this $S_n$... I mean, I can't somehow say enough how much you should make sure you understand $S_n$. Understand the cycle notation, etc. It's extremely important. If you learn anything in this class that you'll learn later in your life, you'll regret not thinking about $S_n$.

So, let me say some more comments about it. We said that... so, maybe I say it again, just because I don't know whether this is all things you've studied before or not, but...

So, $S_n$ I write now \inlinemetanote{writes: $S_n$} --- this is not $S_n$, this is the set of $n$ points. So, I should maybe write it differently \inlinemetanote{erases doodle}. I'll try to say one more thing about the cycle notation, because, as I said, it's extremely important, and it's not like I gave you a totally formal presentation of cycle notation. I gave you a little bit, and it's kind of your responsibility to work out the cycle notation from that. But I'll give you one more hint.
\end{speech}

\begin{speech}[00:29:27 - 00:30:52]
So, if I take here the set $[n]$, which is the set $\{1, 2, 3, \dots, n\}$ \inlinemetanote{writes: $[n] = \{1, 2, 3, \dots, n\}$}, and now I write them as dots... Okay, let's make it $7$. Is that $7$? $4, 5, 6$... Ah, very good \inlinemetanote{draws $7$ points}.

Okay, so, if this is my set, if I take an element in here, in $S_7$ \inlinemetanote{writes: $S_7 \ni \sigma$}, then it maps this set \inlinemetanote{writes: $\sigma \colon [n] \to [n]$}; it's a bijection, that's by definition, it's a bijection from this set to itself.

So, how do I see the cycles if I want to see it as a picture? I pick $1$, and $\sigma$ takes it to somewhere. So, maybe $\sigma$ takes $1$ to $4$ \inlinemetanote{draws an arrow from $1$ to $4$ with $\sigma$}. It has to go somewhere. And then maybe it goes to $6$ \inlinemetanote{draws an arrow from $4$ to $6$}. It keeps going. But eventually this path has to repeat, right? Because there's only finitely many elements. It can't keep going to different points; eventually it has to repeat.
\end{speech}

\begin{speech}[00:30:52 - 00:32:32]
Could it repeat like this? \inlinemetanote{gestures to a branching/intersecting path} Could it be like going here, here, here, here? The answer is no! Because it's a bijection. We already know who comes to that; there can't be something else that comes to that. So, it cannot repeat like that. It has to repeat, and the only way it can repeat is to go back to the beginning. That's the only way it can repeat.

So, another way to think about these cycles is just you look at it dynamically: there's the seven points, or however many points, you pick one, it goes, and eventually it comes back to $1$. And that's the cycle for $1$. In this case, it would be $(1\ 4\ 6)$ \inlinemetanote{writes: $(1\,4\,6)$}.

And then we take this stuff out, and we go with the points that are remaining. Again, now we start with $2$. $2$ has to go, and has to be a cycle, and has to come back to $2$. And it can never interfere with the cycle here because of the bijection property. So, you just take out another cycle. You do that, and that's how you get the cycle decomposition.

Associated to the cycle decomposition, let's pretend here $2$ goes to $3$ and comes back, then this would be $(2\ 3)$ \inlinemetanote{writes: $(2\,3)$}. And what's left? $7$ by himself, all right? That would be in this case. But that's how you visually see it.

You notice --- and this will be for $\sigma$ --- out of this cycle decomposition, you get a certain set of numbers: the lengths of the cycles. And you get $7 = \dots$ I don't know, it's not working, we missed one! \inlinemetanote{chuckles} $5$, sorry, here \inlinemetanote{draws cycle connecting $5$ and $7$}. There! You get $7 = 3 + 2 + 2$ \inlinemetanote{writes: $7 = 3 + 2 + 2$}.

Whenever you have... because at the end you have to account for all the points. Every point has to be in one of the cycles, so you get all the points, and that's the lengths of the cycles.

So, this gives you a connection to partitions, and that's an extremely deep connection. So, this is a partition. Do you know that? That's like an object: how to take a number and break it up into other numbers. There's a very deep connection between the symmetric group and the study of partitions \inlinemetanote{writes: $S_n \longleftrightarrow \text{Partitions}$}.

Maybe I'll say a little bit about this later. But have you ever thought about the study of partitions? Maybe I'll say a little bit about it. But not today. That's a great subject. Okay, but not today.
\end{speech}

\begin{content}[Cycle Decomposition in \texorpdfstring{$S_7$}{S7} and Integer Partitions]
Let $[n] := \{1, 2, \dots, n\}$. Any permutation $\sigma \in S_n$ is a bijection $\sigma \colon [n] \to [n]$.

\definecolor{figBlue}{HTML}{076090} % RGB (7, 96, 144): petrol blue (MidnightBlue)
\definecolor{figRed}{HTML}{B5241C} % RGB (181, 36, 28): brick red (BrickRed)
\definecolor{figGreen}{HTML}{3C8252} % RGB (60, 130, 82): dark green (ForestGreen!70!black)
\definecolor{figGray}{HTML}{666666} % RGB (102, 102, 102): dark gray
\begin{hidden}
- Canvas: no scale, cm; 7 discrete points arranged horizontally at y = 0 from x = 1.0 to 7.0, step 1.0;
  curved arrows showing the 3 disjoint cycles; footprint x in [0.5, 7.8], y in [-1.5, 1.8]
- Order: (1) 7 points with numbers below; (2) Cycle (1 4 6): 1 -> 4 (top arc), 4 -> 6 (top arc), 6 -> 1 (bottom return arc);
  (3) Cycle (2 3): 2 -> 3 (top arc), 3 -> 2 (bottom return arc);
  (4) Cycle (5 7): 5 -> 7 (top arc), 7 -> 5 (bottom return arc);
  (5) Cycle notation and partition equation below.
- Coordinates:
  Points: P_k = (k, 0) for k = 1..7.
  Cycle 1 (figBlue): 1 -> 4 to[bend left=45], 4 -> 6 to[bend left=45], 6 -> 1 to[bend left=35] (below, so from 6 to 1 below is to[bend left=35])
  Cycle 2 (figGreen): 2 -> 3 to[bend left=55], 3 -> 2 to[bend left=55] (below)
  Cycle 3 (figRed): 5 -> 7 to[bend left=45], 7 -> 5 to[bend left=45] (below)
- Labels:
  Numbers 1..7 anchor north at (k, -0.15)
  Cycle label \sigma = (1 4 6)(2 3)(5 7) below at (4.0, -1.1)
  Partition 7 = 3 + 2 + 2 below at (4.0, -1.6)
\end{hidden}
\begin{center}
\begin{tikzpicture}[>=Stealth]
    % (1) Points 1 to 7
    \foreach \i in {1, ..., 7} {
        \coordinate (P\i) at (\i, 0);
        \fill (P\i) circle (2.2pt);
        \node[anchor=north] at (\i, -0.15) {$\i$};
    }

    % (2) Cycle (1 4 6) in Blue
    \draw[->, thick, figBlue] (P1) to[bend left=40] node[midway, above, font=\footnotesize] {$\sigma$} (P4);
    \draw[->, thick, figBlue] (P4) to[bend left=40] (P6);
    \draw[->, thick, figBlue] (P6) to[bend left=35] (P1);

    % (3) Cycle (2 3) in Green
    \draw[->, thick, figGreen] (P2) to[bend left=55] (P3);
    \draw[->, thick, figGreen] (P3) to[bend left=55] (P2);

    % (4) Cycle (5 7) in Red
    \draw[->, thick, figRed] (P5) to[bend left=45] (P7);
    \draw[->, thick, figRed] (P7) to[bend left=45] (P5);
\end{tikzpicture}
\end{center}

In cycle notation, this permutation is written as a product of disjoint cycles:
\[
\sigma = (1\ 4\ 6)(2\ 3)(5\ 7) \in S_7.
\]
The lengths of the disjoint cycles partition the total number of elements:
\[
7 = 3 + 2 + 2.
\]

\begin{steps}[Connection between Symmetric Groups and Integer Partitions]
Every permutation $\sigma \in S_n$ decomposes uniquely into a product of disjoint cycles. The multiset of cycle lengths forms an \newterm{integer partition} of $n$, denoted:
\[
n = \lambda_1 + \lambda_2 + \dots + \lambda_r, \quad \lambda_1 \ge \lambda_2 \ge \dots \ge \lambda_r \ge 1.
\]
Going beyond the lecture: Two permutations in $S_n$ are conjugate if and only if they possess the identical cycle type. Therefore, the conjugacy classes of the symmetric group $S_n$ are in bijective correspondence with the integer partitions of $n$:
\[
S_n / \text{conjugacy} \longleftrightarrow \text{Partitions}(n).
\]
\end{steps}
\end{content}

\subsection{Rotational Symmetries of the Platonic Solids}

\begin{speech}[00:32:32 - 00:34:15]
So, I want to do some examples. And the examples should be symmetry groups of geometric objects. And the most famous geometric objects are the Platonic solids, in some sense the most famous.

So, you could take the tetrahedron... \inlinemetanote{draws a tetrahedron} the cube... \inlinemetanote{draws a cube} and then there's the octahedron... \inlinemetanote{draws an octahedron} I'm not gifted, but that's not terrible. And then there's the dodecahedron, which is much harder to draw. You know, I don't know. Well, you know what this thing is, right? But that's not bad. So, this is the one with $12$ sides, and then of course there's the icosahedron \inlinemetanote{draws an icosahedron}, which is even more complicated to draw.

Those are the five Platonic solids. And one thing in this class, at some point, you should be familiar with what the groups of rotational symmetries of these objects are.

And if you remember a long time ago --- well, couldn't be that long, since the beginning of the class --- we discussed this. For every one of these, there's a group of rotational symmetries, which means you imagine this object, and you look at the rotations that move the faces to the faces.

And we discussed this once before. So, this one, let's just review this: how many rotational symmetries does this have? \inlinemetanote{points to the cube, writes: $\text{how many rotat. symm.}$ under it}
\end{speech}

\begin{speech}[00:34:15 - 00:35:37]
You just imagine you have a cube, and it's really you could just pick up a cube and do this: how many ways can you rotate it? And the group is rotation and composition. The neutral element is do nothing, and the inverse element is undo what you did. That's the proof it's a group.

So, how many rotational symmetries does this have? You can just rotate the cube in your head.

The standard way to count this --- there's many different ways to count it, but one way to count it is you take the top side, and suppose just take the top of the box. You can ask: how many ways can I rotate it, keeping that the top of the box? And the answer is clearly four. If you want to keep it to be the top, there's four things you can do. Do you agree? So, that's four \inlinemetanote{writes: $4$ under the cube}.

And what else can you do? I can move the top of the box to any other side. And how many sides are there? Six. And it turns out that's all you can do, because you can just think about it. You should think about this yourself. That's the only thing you can do: you can either rotate the box or move it to the other side \inlinemetanote{writes: $4 \cdot 6 = 24$ under the cube}.

But these things, as I said, this is the geometric part of the class, and geometry is all about the camera in your head, and at some point, if you want to understand geometry, you have to develop the camera in your head. There's no other way in some sense. Okay, and I can't really help you do that.

So, what about this one? How many rotational symmetries does this have? \inlinemetanote{points to the tetrahedron} This time you should use the same logic.
\end{speech}

\begin{student}[Student Answer]
$12$
\end{student}

\begin{speech}[00:35:38 - 00:35:49]
Yeah, $12$. $12$ is good. Why do you say it's $12$?
\end{speech}

\begin{student}[Student Answer, continued]
If you fix the bottom side, you can rotate it in three ways, and then there are four faces.
\end{student}

\begin{speech}[00:35:50 - 00:36:52]
Yeah, three ways, exactly. And then four faces, yeah, in this way. So, there's lots of different ways, there's no unique way to do it. This one is: $3 \times 4 = 12$ \inlinemetanote{writes: $3 \cdot 4 = 12$ under the tetrahedron}.

And what about this one? \inlinemetanote{points to the octahedron} So, let's just do this together. We could pick one of the triangles... You know what these things are, right? You've seen them, it's not...

So, if you take one of the triangle sides, there's three motions around, and eight ways it can go, so, it's $24$ also \inlinemetanote{writes: $3 \cdot 8 = 24$ under the octahedron}.

We leave these... these are more fun, we'll discuss them later \inlinemetanote{points to the dodecahedron and icosahedron}. These are easy to think about, relatively. This one is a super fun thing to think about.

So, that's a group of order $24$.
\end{speech}

\begin{content}[The Five Platonic Solids and Their Rotational Symmetry Groups]
A \newterm{Platonic solid} is a regular, convex polyhedron in three-dimensional space $\mathbb{R}^3$ whose faces are congruent regular polygons and whose vertices have identical valency.

\definecolor{figBlue}{HTML}{076090} % RGB (7, 96, 144): petrol blue (MidnightBlue)
\definecolor{figGray}{HTML}{666666} % RGB (102, 102, 102): dark gray
\definecolor{figLightGray}{HTML}{A6A6A5} % RGB (166, 166, 165): light gray (gray!70)
\begin{hidden}
- Canvas: no scale, cm; 5 wireframe sketches side-by-side in a horizontal row;
  footprint x in [0.0, 14.5], y in [-1.5, 2.5]
- Spacing: 5 items horizontally, centers at x = 1.3, 4.0, 6.8, 9.8, 12.8; y center around 0.8
- Shapes:
  1. Tetrahedron (center x = 1.3): front triangle (0.3, 0.0) -- (2.3, 0.0) -- (1.3, 1.8) -- cycle;
     inner vertex (1.3, 0.6) connected to all 3 corners.
  2. Cube (center x = 4.0): front square (3.3, 0.0) rectangle (4.5, 1.2);
     back square offset by (0.4, 0.4), dashed hidden edges.
  3. Octahedron (center x = 6.8): upper pyramid apex (6.8, 1.9), lower pyramid apex (6.8, -0.3);
     equatorial square (5.9, 0.8) -- (6.5, 0.5) -- (7.7, 0.8) -- (7.1, 1.1) -- cycle.
  4. Dodecahedron (center x = 9.8): pentagonal profile.
  5. Icosahedron (center x = 12.8): triangular facets.
- Labels below:
  Tetrahedron: "3 x 4 = 12" at y = -0.5
  Cube: "4 x 6 = 24" at y = -0.5
  Octahedron: "3 x 8 = 24" at y = -0.5
  Dodecahedron: "12 faces" at y = -0.5
  Icosahedron: "20 faces" at y = -0.5
- Collisions: ample room between shapes; text below paths.
\end{hidden}
\begin{center}
\begin{tikzpicture}[>=Stealth, thick]
    % --- 1. Tetrahedron ---
    \begin{scope}[xshift=0.3cm]
        \coordinate (T1) at (0, 0);
        \coordinate (T2) at (2.0, 0);
        \coordinate (T3) at (1.0, 1.7);
        \coordinate (T0) at (1.0, 0.6);
        \draw[figBlue] (T1) -- (T2) -- (T3) -- cycle;
        \draw[figBlue] (T1) -- (T0) (T2) -- (T0) (T3) -- (T0);
        \node[anchor=north, align=center, font=\small] at (1.0, -0.2) {\textbf{Tetrahedron}\\$3 \cdot 4 = 12$};
    \end{scope}

    % --- 2. Cube ---
    \begin{scope}[xshift=3.2cm]
        \coordinate (C1) at (0, 0);
        \coordinate (C2) at (1.4, 0);
        \coordinate (C3) at (1.4, 1.4);
        \coordinate (C4) at (0, 1.4);
        \coordinate (C5) at (0.5, 0.5);
        \coordinate (C6) at (1.9, 0.5);
        \coordinate (C7) at (1.9, 1.9);
        \coordinate (C8) at (0.5, 1.9);
        \draw[figBlue] (C1) -- (C2) -- (C3) -- (C4) -- cycle;
        \draw[figBlue] (C2) -- (C6) -- (C7) -- (C3) (C4) -- (C8) -- (C7);
        \draw[figBlue, dashed] (C1) -- (C5) -- (C6) (C5) -- (C8);
        \node[anchor=north, align=center, font=\small] at (0.95, -0.2) {\textbf{Cube}\\$4 \cdot 6 = 24$};
    \end{scope}

    % --- 3. Octahedron ---
    \begin{scope}[xshift=6.4cm]
        \coordinate (Otop) at (1.0, 1.9);
        \coordinate (Obot) at (1.0, -0.1);
        \coordinate (O1) at (0.0, 0.9);
        \coordinate (O2) at (0.7, 0.5);
        \coordinate (O3) at (2.0, 0.9);
        \coordinate (O4) at (1.3, 1.3);
        \draw[figBlue] (Otop) -- (O1) -- (Obot) -- (O2) -- (Otop) -- (O3) -- (Obot);
        \draw[figBlue] (O1) -- (O2) -- (O3);
        \draw[figBlue, dashed] (Otop) -- (O4) -- (Obot) (O1) -- (O4) -- (O3);
        \node[anchor=north, align=center, font=\small] at (1.0, -0.2) {\textbf{Octahedron}\\$3 \cdot 8 = 24$};
    \end{scope}

    % --- 4. Dodecahedron ---
    \begin{scope}[xshift=9.5cm]
        \draw[figBlue] (1.0, 0.9) circle (0.9cm);
        \draw[figBlue] (1.0, 0.9) ++(90:0.5cm) \foreach \a in {162, 234, 306, 18} { -- ++(\a:0.5cm) } -- cycle;
        \node[anchor=north, align=center, font=\small] at (1.0, -0.2) {\textbf{Dodecahedron}\\$5 \cdot 12 = 60$};
    \end{scope}

    % --- 5. Icosahedron ---
    \begin{scope}[xshift=12.5cm]
        \draw[figBlue] (1.0, 0.9) circle (0.9cm);
        \draw[figBlue] (0.5, 0.4) -- (1.5, 0.4) -- (1.0, 1.5) -- cycle;
        \node[anchor=north, align=center, font=\small] at (1.0, -0.2) {\textbf{Icosahedron}\\$3 \cdot 20 = 60$};
    \end{scope}
\end{tikzpicture}
\end{center}

\begin{steps}[The Orbit-Stabilizer Counting Logic for Rotational Symmetries]
For any polyhedron $P \subset \mathbb{R}^3$, its rotational symmetry group $\operatorname{Rot}(P) \le \mathrm{SO}(3)$ acts transitively on the set of faces. The order of the group satisfies:
\[
|\operatorname{Rot}(P)| = (\text{rotational symmetries preserving a chosen face}) \times (\text{total number of faces}).
\]
\begin{itemize}
    \item \textbf{Tetrahedron:} 4 triangular faces. Fixing one face allows 3 in-plane rotations ($120^\circ, 240^\circ, 360^\circ$):
    \[
    |\operatorname{Rot}(\text{Tetrahedron})| = 3 \times 4 = 12.
    \]
    \item \textbf{Cube:} 6 square faces. Fixing one face allows 4 in-plane rotations ($90^\circ, 180^\circ, 270^\circ, 360^\circ$):
    \[
    |\operatorname{Rot}(\text{Cube})| = 4 \times 6 = 24.
    \]
    \item \textbf{Octahedron:} 8 triangular faces (dual to the cube). Fixing one face allows 3 rotations:
    \[
    |\operatorname{Rot}(\text{Octahedron})| = 3 \times 8 = 24.
    \]
    \item \textbf{Dodecahedron:} 12 pentagonal faces. Fixing one face allows 5 rotations:
    \[
    |\operatorname{Rot}(\text{Dodecahedron})| = 5 \times 12 = 60.
    \]
    \item \textbf{Icosahedron:} 20 triangular faces (dual to the dodecahedron). Fixing one face allows 3 rotations:
    \[
    |\operatorname{Rot}(\text{Icosahedron})| = 3 \times 20 = 60.
    \]
\end{itemize}
\end{steps}
\end{content}

\subsection{The Rotational Symmetry Group of the Cube and \texorpdfstring{$S_4$}{S4}}

\begin{speech}[00:36:52 - 00:38:42]
So, these rotations of this cube is a group of order $24$ \inlinemetanote{writes: $\operatorname{Rot}(\square) \text{ is a group of order } 24$}.

What group is it? Do you know any groups of order $24$? Well, $24$ happens to be a factorial: $24$ happens to be $4$ factorial. So, that's a group of order $24$, $4$ factorial. We know another group of order $24$, is $S_4$. $S_4$ is also a group of order $24$.

So, this is now a beautiful geometric question: is this rotational group, these rotations here of this cube --- this is bad terminology, but anyway, it's not great notation, but it's very convincing if I do it: is rotations of the... is it isomorphic to $S_4$? \inlinemetanote{writes: $\text{Is } \operatorname{Rot}(\square) \cong S_4\text{?}$} That's a beautiful question!

And I don't know whether I should solve it for you or not. Maybe I solve it for you; I do this one for you.

There's, again, many ways you can do this. One way to do it is you write down the multiplication table for this rotation group, you name all $24$ elements --- you have to make up your own, now you can be like Fred, Bob, Jane, whatever you want, you name $24$ elements here, you name the $24$ elements for that $S_4$, and those are already named by the cycle decompositions --- and then you find some bijection that respects the product.

That is a hard way to do it. But you could do it that way. It's a hard way to do it because you have two sets of $24$ elements, and very few of them are the correct bijection. So, that's not a great way of doing it.

There is a beautiful way of doing it, which I will show you in a moment. And it's good to show this because it shows, you know, the power of geometry here. \inlinemetanote{erases the middle blackboard}
\end{speech}

\begin{speech}[00:38:42 - 00:40:22]
So, it helps here to have a professor who's artistically gifted, and you do not have this luck, unfortunately. \inlinemetanote{laughter} This is not a hard diagram, I should be able to handle it.

Here is the... there's the cube, all right? \inlinemetanote{draws a 3D wireframe cube} And it turns out there is a set of four things here. There's a beautiful geometric set of four things, and those are the four long diagonals.

So, if I take this point of the cube, there's an opposite point, and there's a long diagonal \inlinemetanote{draws first long diagonal through the center}. Okay? If you imagine a box, for every point in the cube, there's an opposite point. Meaning if that point has coordinates $(+1, +1, +1)$ in the cube, it's $(-1, -1, -1)$ if the cube has side length $2$. So, it has an opposite, and there's one of them.

And there's eight vertices, so there's four long diagonals. So, I'll try to draw them, I don't know if it's so helpful. It's not so good, they all pass through the origin. I think I missed one, too. One, two... where is it? One, two, three... this one... No, this one's hidden, yeah, it's there \inlinemetanote{draws all four diagonals intersecting at origin}. Okay, I told you.

So, there's $1, 2, 3, 4$. But still, as I said, it's much better to have the camera in your head. There's the box, and there's four long diagonals; they all go through the middle, this origin. Is that clear?
\end{speech}

\begin{speech}[00:40:22 - 00:41:47]
And now, you label these as, say, $A, B, C, D$. Where is $D$? $D$ \inlinemetanote{labels the four diagonals $A, B, C, D$}. So, there's four of them; you give them different names.

And you can see now, when you rotate this cube, you permute these four long diagonals. That's just a fact. I mean, you move them, the long diagonals go to other long diagonals.

So, that's kind of an amazing thing. That shows that a rotation of the cube gives you a permutation of the long diagonals. And the proof is just that: the proof is that this rotation of this cube... this is a group, and there's a homomorphism to the permutation group of the long diagonals, which is $S_4$ \inlinemetanote{writes: $\operatorname{Rot}(\square) \xrightarrow{\Phi} S_4$}.

And the claim is that the rotations go to permutations of the long diagonals \inlinemetanote{writes: $\Phi \colon \text{rotations} \to \text{permutations of the long diagonals}$}. I say long diagonals all the time here, but I mean, you know, if you pick two points on the same face, then there's kind of a short diagonal that's on the side. But every time you rotate, you permute the long diagonals, and the claim is this is an isomorphism \inlinemetanote{writes: $\text{This is an isom.}$}.
\end{speech}

\begin{speech}[00:41:47 - 00:42:34]
So, how are you going to prove this is an isomorphism? It's not going to be too hard. You just have to check every time you rotate, you do move some long diagonal; it's not trivial. Because both groups are already of order $24$.

So, if you prove that it's injective, then it has to be surjective. Or you could also prove it's surjective. But you have to think about how when you move this cube, it cannot be the case that if you do a non-trivial rotation, all the long diagonals stay fixed.

So, I leave that to you to think about. But that's a totally geometric proof, a beautiful geometric proof, which is way better than trying to write down this table, which, you know, probably you will not succeed in doing.
\end{speech}

\begin{content}
\begin{theorem}\label[theorem]{thm:cube-rotations-isomorphism}
The group of rotational symmetries of a cube is isomorphic to the symmetric group on four letters:
\[
\operatorname{Rot}(\text{Cube}) \cong S_4.
\]
\end{theorem}

\definecolor{figBlue}{HTML}{076090} % RGB (7, 96, 144): petrol blue (MidnightBlue)
\definecolor{figRed}{HTML}{B5241C} % RGB (181, 36, 28): brick red (BrickRed)
\definecolor{figGreen}{HTML}{3C8252} % RGB (60, 130, 82): dark green (ForestGreen!70!black)
\definecolor{figOrange}{HTML}{CC7F2B} % RGB (204, 127, 43): ochre orange (BurntOrange!80!black)
\definecolor{figPurple}{HTML}{8D1990} % RGB (141, 25, 144): purple
\definecolor{figLightGray}{HTML}{A6A6A5} % RGB (166, 166, 165): light gray (gray!70)
\begin{hidden}
- Canvas: no scale, cm; 3D perspective cube centered at (0, 0); footprint x in [-3.5, 3.5], y in [-2.5, 2.5]
- Geometry:
  Vertices of cube [-1.5, 1.5]^3 projected in oblique view:
  Vector x: (1, 0), Vector y: (0, 1), Vector z: (-0.6, -0.45)
  Front face: V1=(-1.5, -1.5), V2=(1.5, -1.5), V3=(1.5, 1.5), V4=(-1.5, 1.5)
  Back face: shift by z -> V5=V1+z, V6=V2+z, V7=V3+z, V8=V4+z
  Center O = (0, 0)
- 4 Long Diagonals (connecting antipodal vertices through O):
  Diagonal A (figRed): connects V1 to V7
  Diagonal B (figBlue): connects V2 to V8
  Diagonal C (figGreen): connects V3 to V5
  Diagonal D (figOrange): connects V4 to V6
- Order: (1) Back edges of cube (dashed); (2) 4 body diagonals; (3) Front edges of cube (solid); (4) Center origin O; (5) Labels A, B, C, D
- Coordinates:
  V1 = (-1.5, -1.5), V2 = (1.5, -1.5), V3 = (1.5, 1.5), V4 = (-1.5, 1.5)
  V5 = (-2.2, -2.0), V6 = (0.8, -2.0), V7 = (0.8, 1.0), V8 = (-2.2, 1.0)
  Center O = (-0.35, -0.25)
- Labels:
  Diagonal A: at V7, label "$A$" anchor south west
  Diagonal B: at V8, label "$B$" anchor north west
  Diagonal C: at V3, label "$C$" anchor south west
  Diagonal D: at V4, label "$D$" anchor south east
- Collisions: ample margins around the wireframe; text drawn last.
\end{hidden}
\begin{center}
\begin{tikzpicture}[>=Stealth, line cap=round, line join=round]
    % Vertex coordinates
    \coordinate (V1) at (-1.5, -1.5);
    \coordinate (V2) at (1.5, -1.5);
    \coordinate (V3) at (1.5, 1.5);
    \coordinate (V4) at (-1.5, 1.5);

    \coordinate (V5) at (-2.2, -2.0);
    \coordinate (V6) at (0.8, -2.0);
    \coordinate (V7) at (0.8, 1.0);
    \coordinate (V8) at (-2.2, 1.0);
    \coordinate (O) at (-0.35, -0.25);

    % (1) Back edges of the cube (dashed)
    \draw[thick, figLightGray, dashed] (V5) -- (V6) (V5) -- (V8) (V5) -- (V1);

    % (2) Four Long Body Diagonals
    \draw[thick, figRed] (V5) -- (V3);       % Diagonal C
    \draw[thick, figBlue] (V6) -- (V4);      % Diagonal D
    \draw[thick, figGreen] (V1) -- (V7);     % Diagonal A
    \draw[thick, figOrange] (V2) -- (V8);    % Diagonal B

    % Origin intersection
    \fill[black] (O) circle (2pt);

    % (3) Front edges of the cube (solid)
    \draw[thick, figBlue!80!black] (V1) -- (V2) -- (V3) -- (V4) -- cycle;
    \draw[thick, figBlue!80!black] (V2) -- (V6) -- (V7) -- (V3);
    \draw[thick, figBlue!80!black] (V4) -- (V8) -- (V7);

    % (4) Diagonal Labels
    \node[anchor=south west, figGreen, font=\bfseries] at (V7) {$A$};
    \node[anchor=south east, figOrange, font=\bfseries] at (V8) {$B$};
    \node[anchor=south west, figRed, font=\bfseries] at (V3) {$C$};
    \node[anchor=south east, figBlue, font=\bfseries] at (V4) {$D$};
\end{tikzpicture}
\end{center}

\begin{proof}[Proof of \cref{thm:cube-rotations-isomorphism}]
A cube has 8 vertices, which form 4 antipodal pairs. The line segments connecting each antipodal pair through the origin are the 4 \newterm{long diagonals} (space diagonals), denoted by the set:
\[
\Delta := \{A, B, C, D\}.
\]
Every rotation $g \in \operatorname{Rot}(\text{Cube})$ maps vertices to vertices and preserves antipodal pairs, thereby permuting the set $\Delta$. This defines a group homomorphism:
\[
\Phi \colon \operatorname{Rot}(\text{Cube}) \to \operatorname{Sym}(\Delta) \cong S_4, \quad g \mapsto \sigma_g.
\]
We claim that $\Phi$ is an isomorphism:
\begin{itemize}
    \item \textbf{Injectivity:}
    Suppose $g \in \ker(\Phi)$, meaning $g$ fixes all 4 long diagonals. Since $g$ is a rigid rotation, fixing three linearly independent axes in $\mathbb{R}^3$ forces $g = \operatorname{id}$. Thus, $\ker(\Phi) = \{\operatorname{id}\}$, so $\Phi$ is injective.
    \item \textbf{Surjectivity:}
    We have established that $|\operatorname{Rot}(\text{Cube})| = 24$ and $|S_4| = 4! = 24$. Since $\Phi$ is an injective homomorphism between two finite groups of the same cardinality:
    \[
    |\operatorname{Rot}(\text{Cube})| = |S_4| = 24,
    \]
    $\Phi$ is automatically surjective, hence a group isomorphism. \qedhere
\end{itemize}
\end{proof}
\end{content}

\subsection{The Rotational Symmetry Group of the Tetrahedron and \texorpdfstring{$A_4$}{A4}}

\begin{speech}[00:42:34 - 00:43:53]
Okay, so, we have one more moment here, so, what about this one? \inlinemetanote{draws a tetrahedron}

This has order $12$, right? We said that the rotations of this has order $12$ \inlinemetanote{writes: $\operatorname{Rot}(\triangle) \text{ has order } 12$}. This sits inside... there's a homomorphism here to $S_4$ again! \inlinemetanote{writes: $\operatorname{Rot}(\triangle) \to S_4$} Because any rotation permutes the four vertices. It's a different geometric object. This is a much more sophisticated thing, the long diagonals. This one's kind of obvious: every time I rotate, I permute the four vertices.

And this is also injective. But this is not an isomorphism: this is $24$, this is order $12$, it goes to a subgroup. And that subgroup is a very well-known important subgroup, it's called $A_4$. So, the answer --- I'm not defining this now, but it's just the name --- this thing here is actually isomorphic to what's called $A_4$ \inlinemetanote{writes: $A_4 \cong \operatorname{Rot}(\triangle)$}. We'll talk about it more.

It's not the whole thing. You could say: why didn't we do that here? The rotation group also has a homomorphism to the permutation group of the eight vertices! It does. But that's clearly not an isomorphism, because $8$ factorial is a lot bigger than $24$.

Okay, so, I stop. \inlinemetanote{applause from students; video ends}
\end{speech}

\begin{content}
\begin{theorem}\label[theorem]{thm:tetrahedron-isomorphism}
The group of rotational symmetries of a regular tetrahedron is isomorphic to the alternating group on four letters:
\[
\operatorname{Rot}(\text{Tetrahedron}) \cong A_4.
\]
\end{theorem}

\definecolor{figBlue}{HTML}{076090} % RGB (7, 96, 144): petrol blue (MidnightBlue)
\begin{hidden}
- Canvas: no scale, cm; wireframe tetrahedron centered around (0, 0); footprint x in [-2.5, 2.5], y in [-1.5, 2.0]
- Order: (1) Base triangle vertices V1, V2, V3; (2) Apex V4; (3) Front solid edges, base dashed edge; (4) Labels 1, 2, 3, 4
- Coordinates:
  V1 = (-1.5, -0.6), V2 = (1.5, -0.6), V3 = (0.2, 0.2), V4 = (0.0, 1.6)
- Labels:
  "1" at V1 anchor north east
  "2" at V2 anchor north west
  "3" at V3 anchor south west
  "4" at V4 anchor south
- Collisions: vertices placed with clear margins; text drawn last.
\end{hidden}
\begin{center}
\begin{tikzpicture}[>=Stealth, thick]
    \coordinate (V1) at (-1.6, -0.6);
    \coordinate (V2) at (1.6, -0.6);
    \coordinate (V3) at (0.3, 0.2);
    \coordinate (V4) at (0.0, 1.7);

    % Edges
    \draw[figBlue] (V1) -- (V2) -- (V4) -- cycle;
    \draw[figBlue] (V2) -- (V3) -- (V4);
    \draw[figBlue, dashed] (V1) -- (V3);

    % Vertices
    \fill[black] (V1) circle (2pt) (V2) circle (2pt) (V3) circle (2pt) (V4) circle (2pt);

    % Labels
    \node[anchor=north east] at (V1) {$1$};
    \node[anchor=north west] at (V2) {$2$};
    \node[anchor=south west] at (V3) {$3$};
    \node[anchor=south] at (V4) {$4$};
\end{tikzpicture}
\end{center}

\begin{steps}[Comparison: Cube vs. Tetrahedron Permutations]
\begin{itemize}
    \item \textbf{Tetrahedron Action on 4 Vertices:}
    Any rotation of the tetrahedron permutes its 4 vertices $\{1, 2, 3, 4\}$, yielding an injective homomorphism:
    \[
    \Psi \colon \operatorname{Rot}(\text{Tetrahedron}) \hookrightarrow S_4.
    \]
    Since $|\operatorname{Rot}(\text{Tetrahedron})| = 12$ and $|S_4| = 24$, the image is a subgroup of index $2$, which is precisely the \newterm{alternating group} $A_4$ (the group of even permutations).
    \item \textbf{Why Not Act on the 8 Vertices of the Cube?}
    Permuting the 8 vertices of a cube gives an injective homomorphism $\operatorname{Rot}(\text{Cube}) \hookrightarrow S_8$. However, $|S_8| = 8! = 40\,320 \gg 24$, so the image is a tiny subgroup of $S_8$. The 4 long diagonals provide the exact dimensional match $|S_4| = 24$.
\end{itemize}
\end{steps}
\end{content}

\begin{overview}[Summary of Key Concepts]
\begin{itemize}
    \item \textbf{Standing Focus in Finite Group Theory:} Emphasis on finite, non-Abelian groups alongside prototype extreme families: symmetric groups $S_n$ (order $n!$, non-Abelian for $n \ge 3$) and cyclic groups $\mathbb{Z}_n$ (order $n$, Abelian).
    \item \textbf{Element Orders and Cyclic Subgroups:} For any element $x \in G$ in a finite group, its order $\operatorname{Ord}(x) := \min \{m \ge 1 \mid x^m = e\}$ is well-defined, and the generated subgroup $\langle x \rangle = \{e, x, \dots, x^{\operatorname{Ord}(x)-1}\}$ satisfies $|\langle x \rangle| = \operatorname{Ord}(x)$ with $\langle x \rangle \cong (\mathbb{Z}_{\operatorname{Ord}(x)}, +, [0])$.
    \item \textbf{Lagrange's Theorem:} For any finite group $G$ and subgroup $H \le G$, the order of $H$ divides the order of $G$, i.e., $|H| \mid |G|$. Consequently, $\operatorname{Ord}(x) \mid |G|$ for all $x \in G$, and every group of prime order $p$ is cyclic ($G \cong \mathbb{Z}_p$) with no non-trivial proper subgroups.
    \item \textbf{Coset Equivalence and Proof Architecture:} The relation $x \sim y \iff \exists h \in H \colon xh = y$ defines an equivalence relation whose equivalence classes partition $G$. For any class $E$ and representative $x \in E$, the map $\phi \colon H \to E$, $\phi(h) = xh$, is a bijection (injective by group cancellation, surjective by definition of $\sim$). Thus, $|E| = |H|$, yielding $|G| = [G : H] \cdot |H|$.
    \item \textbf{Cycle Decompositions and Partitions:} Every permutation in $S_n$ uniquely decomposes into disjoint cycles. The multiset of cycle lengths forms an integer partition of $n$, connecting the conjugacy structure of $S_n$ to additive combinatorics.
    \item \textbf{Rotational Symmetries of Platonic Solids:}
    \begin{itemize}
        \item \textbf{Cube:} Order $24 = 4 \times 6$. The rigid rotations permute the 4 long space diagonals, establishing the isomorphism $\operatorname{Rot}(\text{Cube}) \cong S_4$.
        \item \textbf{Tetrahedron:} Order $12 = 3 \times 4$. Rotations permute the 4 vertices as even permutations, establishing the isomorphism $\operatorname{Rot}(\text{Tetrahedron}) \cong A_4 \le S_4$.
        \item \textbf{Octahedron:} Dual to the cube with order $24 = 3 \times 8$.
    \end{itemize}
\end{itemize}
\end{overview}